\medskip

\begin{enumerate}
	\item On note par 
	\begin{itemize}
		\item $B$ l'évènement : \og on tire un jeton de couleur blanche \fg{} ;
		\item $N$ l'évènement : \og on tire un jeton de couleur noire \fg{}.
	\end{itemize}
	Ainsi, $p(B)=\dfrac35$ et $p(N)=\dfrac25$
	\begin{enumerate}
			\item L'arbre complété :
\begin{center}
			\pstree[treemode=R,nodesepA=0pt,nodesepB=2.5pt,treesep=1cm,levelsep=2.5cm]{\TR{}}
				{\pstree{\TR{$B$~}\taput{$\frac35$}}
					{\TR{$B$}\taput{$\frac35$}
					\TR{$N$}\tbput{$\frac25$}
					}
				\pstree{\TR{$N$~}\tbput{$\frac25$}}
					{\TR{$B$}\taput{$\frac35$}
					\TR{$N$}\tbput{$\frac25$}
					}
				}
\end{center}
			\item Pour perdre 9 euros, il faut avoir tiré deux jetons blancs.
			
Sachant qu'il s'agit de deux tirages successifs avec remise, et en notant $X$ la variable aléatoire donnant le gain du joueur, on a : $p(X = - 9) = \left(\dfrac35\right)^2=\dfrac{9}{25}$
	\end{enumerate}
\item 
\begin{enumerate}
	\item On considère à présent qu'il y a $3$ jetons blancs et $N$ jetons noirs, pour un total de $N+3$ jetons. On effectue toujours deux tirages avec remise d'un jeton. On note par $X$ la variable aléatoire donnant le gain du joueur. Donc $X\in\lbrace-9~;~-1~;~5\rbrace$
	
	\begin{itemize}
		\item Pour perdre 9 euros, il faut tirer deux jetons blancs : $P(X=-9)=\left(\dfrac{3}{N+3}\right)^2=\dfrac{9}{(N + 3)^2}$
		\item Pour perdre un euro il faut tirer deux jetons noirs : $P(X=-1)=\left(\dfrac{N}{N+3}\right)^2=\dfrac{N^2}{(N + 3)^2}$
		\item pour gagner cinq euros, il faut tirer un jeton de chaque couleur (B + N ou N + B) :\\ $P(X=5)=2\times \dfrac{3}{N + 3} \times \dfrac{N}{N + 3}=\dfrac{6N}{(N + 3)^2}$
	\end{itemize}
La loi de probabilités est donc :
\begin{center}
\begin{tabular}{|l|c|c|c|}\hline
$X=x_i$&$- 9$&$-1$&5 \\ \hline
$P(X=x_i)$&$\dfrac{9}{(N+3)^2}$&$\dfrac{N^2}{(N+3)^2}$&$\dfrac{6N}{(N+3)^2}$ \\ \hline
\end{tabular}
\end{center}
\item $-x^2 + 30x - 81 > 0 \iff x^2-30x+81<0$. Les racines du trinôme $x^2-30x+81$ sont $x=3$ et $x=27$. Donc en utilisant les signes du trinôme du second degré, $-x^2+30x-81>0 \iff x\in ]3~;~27[$.
\item Calculons l'espérance de cette loi de probabilité :

 $E(X)=(-9)\times \dfrac{9}{(N+3)^2} + (-1)\times \dfrac{N^2}{(N+3)^2}+5 \times \dfrac{6N}{(N+3)^2}=\dfrac{-N^2+30N-81}{(N+3)^2}$.

Le jeu est favorable au joueur si et seulement si $E(X)>0$.

$E(X)>0 \iff -N^2+30N-81>0 \iff N\in ]3~;~27[$ donc $4\leqslant N \leqslant 26$ ($N \in \llbracket4~;~26 \rrbracket$) : il faut donc entre 4 et 26 jetons noirs (4 et 26 compris) pour que le jeu soit favorable au joueur.
\item Le gain moyen (ou l'espérance) est maximal lorsque la fonction $f~:~N\mapsto \dfrac{-N^2 + 30N - 81}{(N + 3)^2}$ est maximale, c'est-a-dire lorsque $f'(N) = 0$.

La fonction $f$ est continue et dérivable sur $[0~;~+\infty[$. Pour $N\in [0~;~+\infty[$ :

 $f'(N)=\dfrac{(-2N + 30)\times (N + 3)^2-2(N + 3)\times (-N^2 + 30N - 81)}{\left((N+3)^2\right)^2}=$

$\dfrac{(-2N + 30)(N + 3)-2(-N^2 + 30N - 81)}{(N+3)^3}
 =\dfrac{-2N^2 - 6N + 30N + 90 + 2N^2 - 60N + 162}{(N + 3)^3}=$

$\dfrac{-36N + 252}{(N + 3)^2}=\dfrac{-36(N - 7)}{(N + 3)^2}$

 $f'(N)=0 \iff N=7$.

 Le gain moyen du joueur est maximal lorsque dans l'urne il y a 7 jetons noirs et 3 jetons blancs.
\end{enumerate}
\item Dans l'urne il y a 7 jetons noirs ($N=7$) et 3 jetons blancs, pour un total de 10 jetons. La probabilité de gagner 5 euros est égale à $p=\dfrac{6\times7}{10^2}=\dfrac{21}{50}=0,42$.

On répète ce jeu pour 10 personnes. On note par $Y$ la variable aléatoire donnant le nombre de personnes gagnant 5 euros. Ainsi $Y \in \llbracket0~;~10 \rrbracket$. $Y$ suit une loi binomiale (expérience de Bernoulli n'ayant que deux issues et se répétant indépendamment et à l'identique $n$ fois) de paramètres $n=10$ et $p=\dfrac{21}{50}$ : $Y \sim \mathcal{B}\left(10~;~\dfrac{21}{50}\right)$

$P(Y\geqslant0)=1-P(Y<1)=1-P(Y=0)=1-\left(\dfrac{29}{50}\right)^{10}\approx 0,996$

Donc la probabilité d'avoir au moins un gagnant à 5 euros est d'environ 0,996. C'est donc un évènement (presque) certain.
\end{enumerate}
